% GAL1897_W04 — source-faithful transcription of physical PDF pages 038–039.
% Primary authority: JP2 leaves L0037–L0038, navigated against the 96-page 1897 PDF.
% Substantive source-level mathematical defects are preserved verbatim and ledgered
% separately in SOURCE_ERRORS.csv and MATHEMATICAL_ERROR_CERTIFICATES.md.

\providecommand{\GalSourcePageStart}[4]{}%
\providecommand{\GalSourcePageBoundary}[4]{\space}%
\providecommand{\GalSourceError}[2]{#1}%
\providecommand{\GalHistoricalFootnote}[1]{\footnote{#1}}%
\providecommand{\GalArticleEndRule}{\par\smallskip\begin{center}\rule{0.18\linewidth}{0.4pt}\end{center}}%

\GalSourcePageStart{038}{L0037}{9}{9}

% ENSEG:W04:PDF038:0001
\begin{center}
{\Large\bfseries NOTES\par}
\vspace{0.35em}
{\scriptsize\bfseries ON\par}
\vspace{0.35em}
{\large\bfseries SOME POINTS IN ANALYSIS\textsuperscript{(*)}.\par}
\vspace{0.75em}
\rule{0.12\linewidth}{0.4pt}
\end{center}

% ENSEG:W04:PDF038:0002
\GalHistoricalFootnote{M.~Gergonne's \emph{Annales de Mathématiques}, volume XXI,
page 182 (1830--1831). Owing to a printing error, it reads
\emph{Galais, student at the École Normale}, instead of \emph{Galois}.
\textsc{(J. Liouville.)}}

\begin{center}
\textbf{§ I. — Proof of a theorem of analysis.}
\end{center}

% ENSEG:W04:PDF038:0003
\noindent\textsc{Theorem.} — \emph{Let $Fx$ and $fx$ be any two given functions;
for all $x$ and $h$ we shall have}
\[
\frac{F(x+h)-Fx}{f(x+h)-fx}=\varphi(k),
\]
\emph{where $\varphi$ is a determinate function and $k$ is a quantity between
$x$ and $x+h$.}\GalSourceErrorMark{W04-SE001}

\medskip
% ENSEG:W04:PDF038:0004
\noindent\emph{Proof.} — Indeed, set
\[
\frac{F(x+h)-Fx}{f(x+h)-fx}=P;
\]
from which we obtain
\[
F(x+h)-Pf(x+h)=Fx-Pfx;
\]
which shows that the function $Fx-Pfx$ does not change when $x$ is replaced by
$x+h$. It follows that, unless it remains constant between these limits---which
could occur only in special cases---this function will have one or more maxima
and minima between $x$ and $x+h$.\GalSourceErrorMark{W04-SE004} Let $k$ be the
value of $x$ corresponding to one of them. We shall plainly have
$k=\psi(P)$, where $\psi$ is a determinate function; hence we must also have
$P=\varphi(k)$, where $\varphi$ is another likewise determinate function;
this proves the theorem.\GalSourceErrorMark{W04-SE003}

% ENSEG:W04:PDF038:0005
From this we may conclude, as a corollary, that the quantity
\[
\mathop{\mathrm{lim.}}\frac{F(x+h)-Fx}{f(x+h)-fx}=\varphi(x),
\]
for $h=0$, is necessarily a function of $x$, which
\GalSourceError{proves, \emph{a priori}, the existence of derived functions.}{W04-SE002}

\GalCriticalNote{W04-SE001}{\textbf{Status: proved missing domain hypothesis.}
The quotient is not defined unless $f(x+h)\ne f(x)$. For example, if $f$ is
identically zero, the denominator vanishes for every $x,h$. The universal
printed wording is retained above. Source: PDF038/L0037. Downstream effect:
imposing $f(x+h)\ne f(x)$ repairs only the quotient's domain and does not
repair the later extremum, inverse-function, or derivative-existence steps.
The bounded post-P13 search found no explicit prior notice securely dated in
the searched corpus. Evidence:
\texttt{W04-SE001} in
\nolinkurl{source/critical_baseline/PRINTED_AND_SOURCE_ERRORS.md} and
\texttt{BIBLIOGRAPHIC\_PRIOR\_NOTICE\_MATRIX.csv}.}

\GalCriticalNote{W04-SE004}{\textbf{Status: proved missing regularity
hypothesis.} A discontinuous function can take equal endpoint values, be
nonconstant, and attain no relevant interior extremum. Continuity on the
closed interval, or another hypothesis ensuring attainment, is required.
Source: PDF038/L0037. Downstream effect: continuity or another attainment
hypothesis repairs only the extremum step; the later inverse inference and
corollary remain independently defective. The bounded post-P13 search found no
explicit prior notice securely dated in the searched corpus. Evidence:
\texttt{W04-SE004} in
\nolinkurl{source/critical_baseline/PRINTED_AND_SOURCE_ERRORS.md} and
\texttt{BIBLIOGRAPHIC\_PRIOR\_NOTICE\_MATRIX.csv}.}

\GalCriticalNote{W04-SE003}{\textbf{Status: proved invalid inference.} From
$k=\psi(P)$ one cannot infer a single-valued $P=\varphi(k)$ unless $\psi$ is
injective or an inverse branch is specified. A noninjective $\psi$ with
$\psi(0)=\psi(1)=0$ gives an immediate counterexample. Source: PDF038/L0037.
Downstream effect: injectivity or a specified inverse branch repairs only
$k=\psi(P)\Rightarrow P=\varphi(k)$ and does not establish the theorem under
the other missing hypotheses. The bounded post-P13 search found no explicit
prior notice securely dated in the searched corpus. Evidence:
\texttt{W04-SE003} in
\nolinkurl{source/critical_baseline/PRINTED_AND_SOURCE_ERRORS.md} and
\texttt{BIBLIOGRAPHIC\_PRIOR\_NOTICE\_MATRIX.csv}.}

\GalCriticalNote{W04-SE002}{\textbf{Status: proved false universal claim.}
For $F(t)=|t|$, $f(t)=t$, and $x=0$, the quotient is $|h|/h$: its right-hand
limit is $1$ and its left-hand limit is $-1$. No local wording correction
repairs the claim; a reconstructed theorem would need explicit hypotheses.
Source: PDF038/L0037. Downstream effect: the stated corollary cannot establish
universal existence of derivatives; a reconstructed result requires explicit
domain and regularity hypotheses. The bounded post-P13 search found no explicit
prior notice securely dated in the searched corpus. Evidence:
\texttt{W04-SE002} in
\nolinkurl{source/critical_baseline/PRINTED_AND_SOURCE_ERRORS.md} and
\texttt{BIBLIOGRAPHIC\_PRIOR\_NOTICE\_MATRIX.csv}.}

\GalSourcePageBoundary{039}{L0038}{10}{10}

\begin{center}
\textbf{§ II. — Radius of curvature of curves in space.}
\end{center}

% ENSEG:W04:PDF039:0006
The radius of curvature of a curve at any one of its points $M$ is the
perpendicular drawn from that point to the intersection of the normal plane at
$M$ with the consecutive normal plane, as is readily verified by geometric
considerations.

% ENSEG:W04:PDF039:0007
Accordingly, let $(x,y,z)$ be a point of the curve. The normal plane at this
point has equation
\[
\begin{array}{@{}r@{\qquad}l@{}}
(N) & (X-x)\dfrac{dx}{ds}+(Y-y)\dfrac{dy}{ds}+(Z-z)\dfrac{dz}{ds}=0,
\end{array}
\]
where $X$, $Y$, and $Z$ denote the running coordinates. The intersection of
this normal plane with the consecutive normal plane is given by this equation
together with the following one:
\[
\begin{array}{@{}r@{\qquad}l@{}}
(I) & (X-x)\dfrac{d.\!\left(\dfrac{dx}{ds}\right)}{ds}
 +(Y-y)\dfrac{d.\!\left(\dfrac{dy}{ds}\right)}{ds}
 +(Z-z)\dfrac{d.\!\left(\dfrac{dz}{ds}\right)}{ds}=1,
\end{array}
\]
since
\[
\left(\frac{dx}{ds}\right)^2+
\left(\frac{dy}{ds}\right)^2+
\left(\frac{dz}{ds}\right)^2=1.
\]
Now it is easy to see that plane (I) is perpendicular to plane (N), because
\[
\frac{dx}{ds}\,d.\!\left(\frac{dx}{ds}\right)+
\frac{dy}{ds}\,d.\!\left(\frac{dy}{ds}\right)+
\frac{dz}{ds}\,d.\!\left(\frac{dz}{ds}\right)=0;
\]
hence the perpendicular drawn from $(x,y,z)$ to the intersection of planes
(N) and (I) is simply the perpendicular drawn from the same point to plane
(I). The radius of curvature is therefore the perpendicular drawn from
$(x,y,z)$ to plane (I). This observation yields, very simply, the known
theorems concerning the radii of curvature of curves in space.

\GalArticleEndRule
